Táto bakalárska práca sa zaoberá riešením vybraných kombinatorických problémov pomocou nástrojov na riešenie problému \acrfull{sat}. Vychádza z poznatku, že viaceré diskrétne úlohy možno opísať spoločným spôsobom: pomocou rozhodovacích premenných, ich prípustných stavov a obmedzení, ktoré musí riešenie spĺňať. Cieľom práce je ukázať, ako sa takto formulované problémy dajú previesť na Boolove formuly, následne na inštancie \acrshort{sat} a riešiť jednotným postupom.

V práci sa ako najdôležitejší matematický aparát používa výroková logika. Pre potreby redukcií sú podstatné najmä pojmy splniteľnosti, logickej ekvivalencie, logického dôsledku a konjunktívnej normálnej formy. Na tomto základe je možné uvažovať o probléme \acrshort{sat} ako o rozhodovacom probléme, pri ktorom sa hľadá také ohodnotenie premenných, ktoré spĺňa zadanú formulu. \Acrshort{sat} má v teoretickej informatike osobitné postavenie aj preto, že ide o prvý známy \Acrshort{np}-úplný problém, ako ukázal Cook v práci \textit{The Complexity of Theorem-Proving Procedures} \cite{Cook1971}.

V kapitole \ref{sec:combinatorics} je zavedený problém splniteľnosti obmedzení ako spoločný rámec pre skúmané kombinatorické úlohy. V tomto modeli je riešením konfigurácia, ktorá spĺňa všetky určené obmedzenia. V práci sú pomocou \acrshort{csp} formulované štyri úlohy: vrcholové farbenie grafov, hranové farbenie grafov, problém nezávislej množiny a problém snarkov. Z tejto skupiny aplikácia priamo podporuje vrcholové a hranové farbenie grafov; nezávislá množina a snarky zostávajú teoretickou formuláciou. Postup redukcie doménových vstupov na \acrshort{sat} je uvedený v podkapitole \ref{sec:combinatorics:reduction}. Priame \acrshort{sat} inštancie aplikácia podporuje ako už pripravené vstupy.

Praktická časť práce (kapitoly \ref{sec:design}, \ref{sec:implementation}) opisuje návrh a implementáciu aplikácie, ktorá umožňuje zadávať vybrané problémy, vytvárať ich \acrshort{sat} kódovanie, spúšťať riešenie a zobraziť výsledok. Návrh systému vychádza z doménovo riadeného návrhu, čistej architektúry, architektúry vertikálnych rezov a architektúry mikroslužieb. Pri návrhu a implementácii sme použili princíp asynchrónneho spracovania úloh. Kapitola \ref{sec:implementation} opisuje aj lokálne spustenie a nasadenie do prostredia Azure.

Testovanie aplikácie je zahrnuté v kapitole \ref{sec:testing}, ktorá opisuje testy, testovacie prostredie a pozorované obmedzenia merania.
